\documentclass[10pt,a4paper]{amsart}
\usepackage[margin=19mm]{geometry}
\usepackage{amsmath,amssymb,mathtools}
\usepackage[scr=rsfs,cal=euler]{mathalfa}
\usepackage{xfrac}
\usepackage[unicode,hidelinks,bookmarks=false]{hyperref}
\newcommand{\ie}{i.e.\ }
\newcommand{\cf}{cf.\ }
\newcommand{\resp}{respectively\ }
\newcommand{\Subsection}{\subsection}
\providecommand{\nobreakdash}{\nobreak\hbox{-}}
\setlength{\emergencystretch}{2em}
\multlinegap0pt
\usepackage{amssymb}
\usepackage{enumerate}
\usepackage{mathtools}
\newtheorem{lemma}{Lemma}
\newtheorem{corollary}{Corollary}
\newtheorem{theorem}{Theorem}
\title{Standard Representations of Crystallographic Coxeter Groups}
\author{Emma Booth (University of Warwick)}
\date{Source: arXiv:2609.04179v1 (CC BY 4.0)}
\begin{document}
\maketitle

\noindent\emph{Source and licence.} Standard Representations of Crystallographic Coxeter Groups. Version arXiv:2609.04179v1 (CC BY 4.0), distributed by arXiv under the Creative Commons Attribution 4.0 International licence, http://creativecommons.org/licenses/by/4.0/, as recorded in the arXiv licence metadata for this exact version and shown on the abstract page https://arxiv.org/abs/2609.04179v1. Full licence notice: \texttt{licenses/arxiv-2609.04179-CC-BY-4.0.txt}.\par\smallskip

\noindent\textit{Editorial scope.} Counted unit: the article's theorem thm:reflectisomorph (source0.tex lines 798--798). The article's complete original proof of it is delivered with it (source0.tex lines 800--818, one \texttt{proof} environment of 19 lines). No mathematics is written, repaired or re-derived: the statement and the proof are the article's own lines, in the article's own notation, and no retained line is substituted or re-flowed.  Retained uncounted as the source-local context the proof needs: lines 759--759; lines 794--796.  Nothing was omitted from any retained span, so the package carries no omission record.  No retained line contains a citation, so the wrapper carries no bibliography and no bibliography entry.  No retained line contains a figure environment, an image-inclusion command or drawing code, so no image is delivered and the package declares no asset dependency.  Editorial formatting only, in addition: an AMS \texttt{amsart} preamble that restates the article's own declarations and macros used by the retained lines, namely the environments \texttt{lemma}, \texttt{corollary}, \texttt{theorem} on separate counters, together with the standard packages those lines need.  The retained environments print in this document as Lemma 1, Corollary 1, Theorem 1 in order of appearance, whereas the article prints the same items with the numbering of its own full document.  The complete native source of the article as distributed in the arXiv e-print for this exact version is bundled unchanged as \texttt{sources/209351/source0.tex}.
\begin{lemma} \label{lem:isomorphcond} Suppose $i \neq j, A_{ij} \neq 0$. Then $\lambda_i = \frac{B_{ij}}{A_{ij}} \lambda_j$, where $B$ is the matrix with entries $B_{ij} = -b_{ij}$.\end{lemma}

\begin{corollary} \label{cor:symmreflectcond} Suppose $A_{ij} A_{ji} \in \left\{0,1,2,3,4\right\}$ for all $i, j$. Then $A$ is symmetrisable if and only if \begin{equation*}
    B_{i_1i_2} \ldots B_{i_k i_1} = A_{i_1i_2} \ldots A_{i_k i_1}
\end{equation*} for all circuits $i_1, \ldots, i_k, i_1$ in the Dynkin diagram. \end{corollary}

\begin{theorem} \label{thm:reflectisomorph} Suppose $A_{ij} A_{ji} \in \left\{0,1,2,3,4\right\}$ for all $i, j$, and $A$ is symmetrisable. Then $H$ and $R$ are isomorphic.\end{theorem}

\begin{proof} By Lemma \ref{lem:isomorphcond}, it is sufficient to show that we can make a well-defined choice of $\lambda_1, \ldots, \lambda_n \in \mathbb{R}$ such that $\lambda_i = \frac{B_{ij}}{A_{ij}} \lambda_j$. 

Fix an arbitrary $i$ and let $\lambda_i = 1$. For each $j$, choose a path $i = k_1, k_2, \ldots, k_r = j$  from $i$ to $j$ and define values of $\lambda_{k_l}$ along this path by \begin{equation*}\lambda_{k_{l+1}} = \dfrac{B_{k_l k_{l+1}}}{A_{k_{l+1} k_l}} \lambda_{k_l}\end{equation*} This will eventually define $\lambda_j$.

Ambiguity in this definition arises where we have two distinct paths $i = k_1, \ldots, k_r = j$ and $i = l_1, \ldots, l_s = j$ from $i$ to $j$. Then we could define either \begin{equation*} \lambda_j = \dfrac{ B_{k_1k_2} \ldots B_{k_r k_1}}{A_{k_1k_2} \ldots A_{k_r k_1}} \lambda_i\end{equation*} or \begin{equation*} \lambda_j = \dfrac{B_{l_1l_2} \ldots B_{l_s l_1}}{A_{l_1l_2} \ldots A_{l_s l_1}} \lambda_i\end{equation*}

For these definitions to be consistent with each other, we need  \begin{equation*}\dfrac{B_{k_1k_2} \ldots B_{k_r k_1}}{A_{k_1k_2} \ldots A_{k_r k_1}}= \frac{B_{l_1l_2} \ldots B_{l_s l_1}}{A_{l_1l_2} \ldots A_{l_s l_1}}\end{equation*}

But if we apply Corollary \ref{cor:symmreflectcond} to the "circuit" $k_1, \ldots, k_r, k_{r-1}, \ldots, k_1$ we see that \begin{equation*}(B_{k_1k_2} \ldots B_{k_r k_1})^2 = A_{k_1k_2} \ldots A_{k_{r-1} k_r} A_{k_2k_1} \ldots A_{k_r k_{r-1}}\end{equation*} and similarly \begin{equation*}(B_{l_1l_2} \ldots B_{l_{s-1} l_s})^2 = A_{l_1l_2} \ldots A_{l_{s-1} l_s} A_{l_2l_1} \ldots A_{l_s l_{s-1}}\end{equation*}

So, squaring both sides and substituting these results, it remains to show that:

\begin{equation*} \dfrac{A_{k_2k_1}\ldots A_{k_rk_{r-1}}}{A_{k_1 k_2} \ldots A_{k_{r-1} k_r}}= \dfrac{A_{l_2l_1} \ldots A_{l_s l_{s-1}}}{A_{l_1 l_2} \ldots A_{l_{s-1} l_s}}\end{equation*}

and multiplying out denominators gives the condition for symmetrisability for the cycle $i = k_1, k_2 \ldots, k_r = j = l_s, l_{s-1}, \ldots, l_2, l_1 = i$. 

So $\lambda_1, \ldots, \lambda_n$ are well-defined. Now let $j, k$ be such that $A_{jk} \neq 0$. Choose a path $l_1, \ldots, l_r$ from $i$ to $j$. Then appending $k$ to this gives a path from $i$ to $k$. So we can use these paths to define $\lambda_j$ and $\lambda_k$, and then substituting in the resulting expression gives \begin{equation*}\lambda_k = \frac{B_{jk}}{A_{kj}} \lambda_j\end{equation*} which is the required condition.

Thus $\varphi: H \to R$ defined by $\varphi(\alpha_i) = \lambda_i e_i$ for each $i$ is the required isomorphism. \end{proof}

\end{document}
